% Insert this helper near the top of main.tex, after the existing \payoff macro.
\newcommand{\pdfplot}[1]{%
  \begin{tikzpicture}[baseline=(current bounding box.center)]
    \begin{axis}[
      width=\linewidth,
      height=4.8cm,
      xmin=0, xmax=50,
      ymin=0, ymax=1.15,
      axis lines=left,
      xtick={25},
      xticklabels={$K$},
      ytick=\empty,
      xlabel={$S$},
      ylabel={pdf},
      clip=false
    ]
      #1
    \end{axis}
  \end{tikzpicture}}

% Insert these slides immediately after:
% \begin{frame}{Options Let You Express Views on the Distribution}
% ...
% \end{frame}

\begin{frame}{Straddle}
  \begin{columns}[T]
    \begin{column}{0.53\textwidth}
      \pdfplot{
        \addplot[very thick, blue, domain=0:50, samples=300]
          {0.80*exp(-((x-25)^2)/72) + 0.05*exp(-((x-25)^2)/260)};
        \addplot[very thick, red!80!black, domain=0:50, samples=300]
          {0.56*exp(-((x-25)^2)/150) + 0.10*exp(-((x-25)^2)/420)};
        \node[anchor=north west, blue] at (rel axis cs:0.04,0.96) {\scriptsize before};
        \node[anchor=north east, red!80!black] at (rel axis cs:0.96,0.96) {\scriptsize more variance};
      }
    \end{column}
    \begin{column}{0.43\textwidth}
      \centering
      \payoff{
        \draw[very thick, sqt-themecolour] (1.0,3.0) -- (3.0,0) -- (5.0,3.0);
        \draw[dashed, gray] (3.0,-0.2) -- (3.0,3.2);
      }
    \end{column}
  \end{columns}

  \vspace{0.2cm}
  \begin{itemize}
    \item Same centre, but more mass moves into the tails away from $K$.
    \item A straddle likes larger moves in either direction, so wider pdf means more value.
  \end{itemize}
\end{frame}
\note{
  This is the clean vol slide. Hold the centre fixed at the strike and spread
  the distribution out. The only thing that changed is variance, and that puts
  more probability on outcomes where $|S-K|$ is large. That is exactly the
  straddle payoff, so long straddles are long variance / long vol.
}

\begin{frame}{Risk Reversal}
  \begin{columns}[T]
    \begin{column}{0.53\textwidth}
      \pdfplot{
        \addplot[very thick, blue, domain=0:50, samples=300]
          {0.72*exp(-((x-25)^2)/100) + 0.07*exp(-((x-25)^2)/320)};
        \addplot[very thick, red!80!black, domain=0:50, samples=300]
          {0.62*exp(-((x-27)^2)/88) + 0.14*exp(-((x-16)^2)/230)};
        \node[anchor=north west, blue] at (rel axis cs:0.04,0.96) {\scriptsize before};
        \node[anchor=south west, red!80!black] at (rel axis cs:0.04,0.08) {\scriptsize more negative skew};
      }
    \end{column}
    \begin{column}{0.43\textwidth}
      \centering
      \payoff{
        \draw[very thick, sqt-themecolour] (0.8,2.6) -- (2.0,1.4) -- (4.0,1.4) -- (5.6,0.0);
        \draw[dashed, gray] (2.0,-0.2) -- (2.0,2.8);
        \draw[dashed, gray] (4.0,-0.2) -- (4.0,2.0);
      }

      \vspace{0.3cm}
      \textbf{Long put / short call risk reversal}
    \end{column}
  \end{columns}

  \vspace{0.2cm}
  \begin{itemize}
    \item More negative skew in the pdf means a fatter left tail.
    \item Downside protection becomes more valuable than upside exposure, which helps a long put / short call structure.
  \end{itemize}
\end{frame}
\note{
  Use this to define skew in plain English: one tail is getting heavier, and
  here it is the downside tail. The risk reversal is the natural skew trade
  because it is long downside convexity and short upside convexity. As skew gets
  more negative, puts should richen relative to calls.
}

\begin{frame}{Butterfly}
  \begin{columns}[T]
    \begin{column}{0.53\textwidth}
      \pdfplot{
        \addplot[very thick, blue, domain=0:50, samples=300]
          {0.70*exp(-((x-25)^2)/70) + 0.05*exp(-((x-15)^2)/180) + 0.05*exp(-((x-35)^2)/180)};
        \node[anchor=north east, blue] at (rel axis cs:0.96,0.96) {\scriptsize underlying pdf};
      }
    \end{column}
    \begin{column}{0.43\textwidth}
      \centering
      \payoff{
        \draw[thick, sqt-themecolour!55] (0.6,0.0) -- (1.8,1.2) -- (3.0,0.0);
        \draw[thick, sqt-themecolour!75] (1.6,0.0) -- (3.0,2.0) -- (4.4,0.0);
        \draw[thick, sqt-themecolour] (3.0,0.0) -- (4.2,1.1) -- (5.4,0.0);
        \draw[dashed, gray] (1.8,-0.2) -- (1.8,1.4);
        \draw[dashed, gray] (3.0,-0.2) -- (3.0,2.2);
        \draw[dashed, gray] (4.2,-0.2) -- (4.2,1.3);
      }

      \vspace{0.3cm}
      \textbf{Stacked butterflies}
    \end{column}
  \end{columns}

  \vspace{0.2cm}
  \begin{itemize}
    \item Overlaying butterflies across strikes gives a piecewise approximation to the underlying pdf.
    \item Each narrow fly is a local density bet around its body strike.
  \end{itemize}
\end{frame}
\note{
  This is the bridge to the more advanced idea. A single narrow butterfly is a
  local bet on density near its middle strike. If you line up lots of them
  across strikes, the weights trace out an approximation to the market-implied
  pdf of the underlying.
}
